\begin{abstract}
Prediction of residential electricity consumption is a challenging time-series problem, primarily because demand varies wildly from 1 household to another based on individual house-to-house usage patterns. A lot of factors like long range monthly and weekly rhythmns as well as short term factors like high power appliances heavily influence the electricity consumption. In this paper, we present a newer, more modern method of electricity demand prediction by framing it time-series problem that can be solved using the transformer and attention mechanism. Using the Temporal Fusion Transformer (TFT) model, along with various engineered power features to improve input feature quality as well as application of preprocessing techniques like Principle Component Analysis (PCA) and Wavelet as well as Bayesian Hyperparameter search, we were able to attain a peak performance on the test set of MAE: 78.57, RMSE: 257.02, R$^2$ 0.8674 and MAPE: 23.57\%. During our experiements, we created 4 variations of the initial model to test the effect of PCA, Wavelet and Bayesian Hyperparameter search on the performance of the model. We also compared our results to the base paper LSTM-SVR study on our dataset \cite{reddy2022stacking} and found that TFT is better suited to residential forecasting because it utilizes the attention mechanism. Further, we also compared it to older baseline architectures like ARIMA and SARIMA. 
\end{abstract}
 
\begin{IEEEkeywords}
Energy Forecasting, Temporal Fusion Transformer, Wavelet Decomposition, Principle Component Analysis, Residential Metering
\end{IEEEkeywords}
 
\section{Introduction}
 
Predicting energy consumption for households is a major part of smart grid optimization for efficient delivery of electricity. Unlike industrial loads which are generally steady and highly predictable, like the Steel Industry Energy consumption dataset used in our base paper, residential consumption data is highly chaotic. It is shaped by the behaviour of the residents, which varies greatly from household to household. 

Older, more traditional methods like ARIMA highly struggle with residential energy predictions since it is not linear nor is it steady and predictable. Looking into newer methods like the Stacking Ensemble appraoch introduced by the base paper has also proved insufficient in this scenario given the highly chaotic nature of the input data. Power consumption varies in a more unpredictable manner as compared to something like an industrial consumption data source. Furthermore, as part of our initial testing, we also experimented with single models like LSTMs which are built to work well with time series data, yet it fails when trying to predict household consumption. 

The Temporal Fusion Transformer (TFT) is highly suited for this kind of multi-horizon energy forecasting. It handles multi meter input as well as variable importance automatically, internally giving relevant attention to the right parameters. 

To improve upon our own initial V1 baseline which had no special preprocessing done, we carried out 3 more experiements to discover what can push the performance further for more accurate predictions. Specifically, we looked at Discrete Wavelet Transform (DWT) which are suited for breaking down an input signal like power usage into various sub-bands, separating long-term trends from short term spikes. Furthermore, we also experimented with the usage of PCA-12 to clean up redundancy among our input data as well as Bayesian Optimization (via Optuna) to find optimal hyperparameters for an efficient model. 

In summary, this paper considers 3 things. First, we built a full TFT piepline for the THC residential dataset. Second, we tested our 4 versions: a baseline TFT V1, TFT + PCA V2, TFT + PCA + Optuna V3, TFT + PCA + Wavelet V4. Finally, we compared and contrasted the results we obtained from each model with the base paper, as well as dove into the specifics of why we got certains types of results. 

\section{Related Work}
 
\subsection{Base Paper: LSTM--SVR Study}

Our base paper compared LSTM variants and SVR on the UCI Steel Industry Energy Consumption dataset \cite{reddy2022stacking}. That study found a large performance gap favoring SVR with an RBF kernel and feature scaling (SVR R$^2$ = 0.9980, best LSTM R$^2$ = 0.5425). The result is important for our current work because it shows that model
behavior is strongly dataset-dependent: methods that work extremely well for stable industrial demand do not automatically indicate that it'll work just as well for more chaotic and noisy residential usage.
 
\subsection{ARIMA and SARIMA Baselines}

ARIMA and SARIMA remain our standard statistical baselines for short-term consumption forecasting.
ARIMA models linear autocorrelation after differencing, while SARIMA goes one step further with explicit seasonal terms. 
These models perform quite well when the input dataset is relatively smooth and seasonal, but in our case of household energy consumption, that is not the case. 
 
\subsection{Temporal Fusion Transformers}

Temporal Fusion Transformers (TFT) were introduced by Lim-et-al. for multi horizon forecasting with mixed covariates. It features a 7 stage hybrid architecture of models, including LSTMs to capture short term dependancies and attention to capture longer term dependancies. This structure matches household usage quite well and thus is a good fit for forecasting demand in households over time. 
 
\section{Dataset and Preprocessing}
 
\subsection{THC Residential Consumption Dataset}
All experiments used a multi meter residential dataset collected at 15-minute level resolution spanning from August 2019 to October 2020 from Uruguay, with eleven monthly CSV files. These were concatenated and aligned on a common datetime index into 1 monolith dataset. Each row contains per-phase active power (\texttt{apower\_ph1}), reactive power, phase voltage, current, power factor, and cumulative energy for one meter reading. After concatenation the dataset contains approximately 144,256 rows across multiple meter identifiers. Missing values imputations were done meter wise using forward-fill followed by backward-fill. Redundant phase-2 and phase-3 energy columns were dropped.
 
\subsection{Feature Engineering}
A base, primary feature set was built identically for all models: apparent power ($S = V \cdot I$), power factor angle ($\theta = \arccos(\text{pf})$), energy intensity, reactive-to-active power ratio ($Q/P$), hour-of-day, day-of-week, binary weekend and peak-hours (08:00--20:00) flags, and a three-step rolling mean of active power.
 
\subsection{Train / Validation / Test Split}
Data is split chronologically per meter to prevent data leakage: 70\% training, up to 85\% validation (cumulative), full series as test. This yields approximately 100,977 , 122,615 and 144,256 samples respectively.